\documentclass[french]{article}
\usepackage[utf8]{inputenc}
\usepackage[T1]{fontenc}
\usepackage{lmodern}
\usepackage{geometry}
\usepackage{graphicx}
\usepackage{babel}
\usepackage{amsmath}
\usepackage{amsthm}
\usepackage{amssymb}
\usepackage{hyperref}
\usepackage{stmaryrd}
\usepackage{enumitem}
\usepackage{tcolorbox}
\usepackage{dsfont}
\usepackage{multirow}
\usepackage{etoc}
\usepackage{titlesec}
\usepackage{esint}
\usepackage{cancel}
\usepackage{mathdots}

\setlist[itemize]{label=$\bullet$}


\renewcommand{\P}{\mathbb{P}}
\newcommand{\N}{\mathbb{N}}
\newcommand{\Z}{\mathbb{Z}}
\newcommand{\Q}{\mathbb{Q}}
\newcommand{\R}{\mathbb{R}}
\newcommand{\C}{\mathbb{C}}
\newcommand{\K}{\mathbb{K}}
\renewcommand{\L}{\mathbb{L}}
\newcommand{\U}{\mathbb{U}}
\newcommand{\st}{^*}
\newcommand{\tq}{\middle|}
\newcommand{\inv}{^{-1}}
\newcommand{\E}{\mathbb{E}}
\newcommand{\F}{\mathbb{F}}
\newcommand{\V}{\mathbb{V}}
\renewcommand{\ker}{\text{Ker}}
\newcommand{\im}{\text{Im}}
\newcommand{\card}{\text{Card}}
\newcommand{\Tr}{\text{Tr}}

\title{TD Euclidiens\\ Exercice 28}
\date{}
\begin{document}
\maketitle

\section{Énoncé}
Soit $E$ un espace euclidien.
\begin{enumerate}
	\item Soit $f\in\mathcal{L}(E)$ tel que $\forall x\in E,\ (f(x)\mid x)=0$. Montrer que $f=0$.
	\item Soit $f\in\mathcal{L}(E)$ autoadjoint tel que $\forall x\in E,\ (f(x)\mid x)=\lVert f(x)\rVert^2$. Montrer que $f$ est un projecteur orthogonal.
\end{enumerate}
\section{Solution}
Remarquer que pour la deuxième question, il suffira de prouver que $f$ est un projecteur car par théorème de cours, un projecteur est orthogonal si, et seulement si, il est autoadjoint.
\begin{enumerate}
	\item $f$ est autoadjoint donc diagonalisable à valeurs propres réelles. Soit $\lambda$ une valeur propre de $f$ et $x$ un vecteur propre. Comme $(f(x)\mid x)=0$, on en déduit que $\lambda\lVert x\rVert^2=0$ donc que $\lambda=0$ car $x$ est non nul en tant que vecteur propre. Donc $f$ est diagonalisable et toutes ses valeurs propres sont nulles donc $f=0$.
	\item On cherche à prouver que $f^2=f$ en se ramenant à la question précédente. Remarquons que $g=f^2-f$ est autoadjoint. Et $$(g(x)\mid x)=(f^2(x)\mid x)-(f(x)\mid x)=(f(x)\mid f(x))-\lVert f(x)\rVert^2=0$$ Donc $g=0$ par la question précédente. Donc $f$ est un projecteur. Donc $f$ est un projecteur orthogonal car il est autoadjoint.
\end{enumerate}
\end{document}
